\chapter{Dijkstra's Algorithm}
\label{ch:ch03}
%% Chapter 3 of the Searchbook.  Code: code/ch03_dijkstra.py (self-test:
%% python3 code/ch03_dijkstra.py).  Figures: figures/ch03/*.tex; the grid,
%% wavefront and expansion figures are generated by code/figures/gen_ch03_*.py.
%% Every number quoted in this chapter is produced and checked by that code.
\SetKwFunction{DijkstraSearch}{Dijkstra}
\SetKwFunction{DijkstraBackward}{BackwardDijkstra}

\begin{objectives}
\begin{itemize}
  \item Explain what Dijkstra's algorithm computes (the cheapest cost from one source to every
vertex of a graph with non-negative edge costs) and when it is the right tool.
  \item Trace it by hand on a small weighted graph and on a grid, keeping track of the priority
queue, the settled set and the parent pointers.
  \item Prove the settled-node invariant, explain why it needs non-negative edge costs, and
derive the running time $\bigO{(\abs{V}+\abs{E})\log\abs{V}}$.
  \item Implement it in \python with \code{heapq} and lazy deletion, on explicit graphs and on
occupancy grids, with the early exit (uniform-cost search) and path reconstruction.
  \item Use one backward run from the goal as an exact heuristic, the table that feeds \astar
(\cref{ch:ch04}) and the single-agent searches inside the multi-agent solvers of
\cref{ch:ch08,ch:ch09}.
\end{itemize}
\end{objectives}

% ---------------------------------------------------------------------
\section{Why this matters for a drone swarm}
\label{sec:ch03-motivation}

Picture one delivery drone at a depot on the edge of a city. Its map is the occupancy grid of
\cref{ch:ch02}: cells are free or blocked, and moving into a cell costs one unit of time, or
more when the cell lies in a wind-exposed corridor or above a schoolyard where the operator has
attached a risk penalty. No-fly zones are simply removed from the grid. The customer is twenty
blocks away. Before the drone lifts off, the planner has to answer one question: which route is
the cheapest, and how do we know that no cheaper route exists?

This is the \textbf{single-source shortest-path problem}\index{shortest path!single-source} on a
graph with non-negative edge costs, and Dijkstra's algorithm \cite{dijkstra1959note} solves it
exactly. It is the oldest algorithm in this book, and the one you will use most often, for three
reasons.
\begin{enumerate}
  \item \emph{It is the baseline planner.} Every other single-agent planner in
\cref{ch:ch04,ch:ch05,ch:ch06} returns a path whose cost is judged against the cost that
Dijkstra's algorithm would have found. When the experiments of \cref{ch:ch25} report a path
length ``relative to the optimum'', the optimum comes from this chapter.
  \item \emph{It is the heuristic factory.} \astar (\cref{ch:ch04}) needs, for every vertex, a
lower bound on the cost that remains to the goal. One run of Dijkstra's algorithm
\emph{backwards} from the goal produces the exact cost-to-go for every vertex of the map. That
table is the perfect heuristic for \astar, and it is the heuristic that the single-agent
searches inside prioritized planning (\cref{ch:ch08}) and conflict-based search (\cref{ch:ch09})
use when they route one drone at a time through space and time.
  \item \emph{It is the ancestor.} \astar is Dijkstra's algorithm with a heuristic added to the
priority; uniform-cost search is Dijkstra's algorithm with an early exit; \dstarlite
(\cref{ch:ch05}) is a backward search from the goal that repairs its distance table when the map
changes. Understand this chapter thoroughly and the next three chapters become variations on one
theme.
\end{enumerate}

In the four-layer hybrid architecture of \cref{ch:ch24}, Dijkstra's algorithm therefore lives in
the global planning layer: it plans a single drone on a static map, and it supplies the distance
tables that the multi-agent solver and the replanning layer consume. It knows nothing about
other drones or about moving intruders; those are the subjects of Parts III and IV.

\paragraph{Roadmap.} \Cref{sec:ch03-intuition} gives the picture of a wave spreading over the
map. \Cref{sec:ch03-problem} fixes the notation, \cref{sec:ch03-algorithm} states the algorithm
with a binary heap and lazy deletion, and \cref{sec:ch03-example} traces it on a seven-vertex
graph and on a $5\times5$ grid. \Cref{sec:ch03-properties} proves that the algorithm is correct,
shows what goes wrong with a negative edge, and derives the running time.
\Cref{sec:ch03-variants} collects the variants you will meet later in the book, above all the
backward run that produces exact heuristics. \Cref{sec:ch03-implementation} discusses the
\python implementation and its pitfalls, and \cref{sec:ch03-drone} places the algorithm in the
drone system.

% ---------------------------------------------------------------------
\section{The idea in plain words}
\label{sec:ch03-intuition}

Drop a stone into a pond. The ripple spreads at constant speed, reaches nearby points before
distant ones, bends around a rock and squeezes through a gap in a wall. If you write down the
moment at which the ripple first reaches each point, you have written down the distance from the
stone to every point of the pond, because the ripple always arrives along a shortest route.
\Cref{fig:ch03-wavefront} shows this picture on a grid map. The bands are the cells whose cost
from the source lies between $k$ and $k+1$, and band $k$ is completely settled before the wave
enters band $k+1$.

\begin{figure}[tb]
  \centering
  \inputfigure{ch03/wavefront}
  \caption{Dijkstra's algorithm as a spreading wave on an 8-connected grid (diagonal moves cost
$\sqrt2$). Cells are shaded by the band $\lfloor\gcost\rfloor$ of their cost from the source;
the wave settles every cell of band $k$ before any cell of band $k+1$, bends around the block
and passes through the gap at the top of the wall. Behind the wall the bands are centred on the
gap, not on the source: cost is measured along paths, not as the crow flies.}
  \label{fig:ch03-wavefront}
\end{figure}

Dijkstra's algorithm\index{Dijkstra's algorithm} simulates this wave on a graph, one vertex at a
time. Every vertex is in one of three states.
\begin{itemize}
  \item \emph{Unseen}: no path to it has been found; its tentative cost is $\gcost(v)=\infty$.
  \item \emph{Open}: at least one path to it is known, with tentative cost $\gcost(v)$, but a
cheaper path may still turn up. The open vertices form the frontier of the wave and are stored
in a priority queue keyed by $\gcost$.
  \item \emph{Settled} (or \emph{closed}): the wave has reached it, and $\gcost(v)$ is final.
\end{itemize}
The loop is the same at every step: take the open vertex with the smallest tentative cost,
declare it settled, and offer each of its neighbours the path that goes through it. If that path
is cheaper than the neighbour's current tentative cost, update the cost and the parent pointer
of the neighbour and put it on the frontier. This update is called
\textbf{relaxation}\index{relaxation}.

\begin{keyidea}
Always settle the open vertex with the smallest tentative cost. Because no edge has a negative
cost, every path that is still to be discovered starts from a vertex that is at least as
expensive, so nothing found later can undercut it: the cost of a vertex is final the moment it
is settled.
\end{keyidea}

Two consequences follow, and both are worth fixing in your mind now. First, vertices are settled
in non-decreasing order of cost; that is the wave. Second, if you only want the cost of one
target $t$, you may stop the moment $t$ is settled: everything settled before $t$ was at least
as cheap, everything still open is at least as expensive, and the answer cannot change any more.
With this early exit the algorithm is what the artificial-intelligence literature calls
\textbf{uniform-cost search}\index{uniform-cost search} \cite{russell2020aima}.

% ---------------------------------------------------------------------
\section{Problem statement and notation}
\label{sec:ch03-problem}

We use the graphs of \cref{ch:ch02}: an explicit graph is a set of vertices with adjacency
lists, and a grid is an implicit graph whose neighbours are generated on demand by a successor
function.

\begin{definition}[Weighted graph, path, path cost]
\label{def:ch03-graph}
A \textbf{weighted directed graph}\index{graph!weighted} is a triple $G=(V,E,c)$ with a finite
vertex set $V$, an edge set $E\subseteq V\times V$ and a cost function $c\colon E\to\R_{\ge0}$.
We write $c(u,v)$ for the cost of the edge $(u,v)$ and $\Succ(u)=\set{v : (u,v)\in E}$ for the
successors of $u$. An undirected edge $\set{u,v}$ of cost $w$ stands for the two directed edges
$(u,v)$ and $(v,u)$, both of cost $w$. A \textbf{path} from $u$ to $v$ is a sequence
$\pi=(v_0,v_1,\dots,v_k)$ with $v_0=u$, $v_k=v$ and $(v_{i-1},v_i)\in E$ for $i=1,\dots,k$; its
cost is $c(\pi)=\sum_{i=1}^{k}c(v_{i-1},v_i)$.
\end{definition}

\begin{definition}[Shortest-path distance]
\label{def:ch03-distance}
The \textbf{shortest-path distance}\index{shortest path!distance} $\delta(u,v)$ is the smallest
cost of any path from $u$ to $v$, and $\delta(u,v)=\infty$ if no such path exists. A path with
$c(\pi)=\delta(u,v)$ is a \textbf{shortest path}; there may be several.
\end{definition}

Three facts about $\delta$ are used throughout the chapter. For all $u,v,w\in V$,
\begin{equation}
\label{eq:ch03-triangle}
\delta(u,u)=0,\qquad
\delta(u,w)\le\delta(u,v)+\delta(v,w),\qquad
\delta(u,v)\le\delta(u,w)+c(w,v)\ \text{ if }(w,v)\in E .
\end{equation}
The second is the triangle inequality: going via $v$ can only be a detour. The third is the
triangle inequality with a single edge as the second leg. A fourth fact is that \emph{subpaths
of shortest paths are shortest paths}: if $\pi$ is a shortest path from $u$ to $v$ and $w$ lies
on $\pi$, then the prefix of $\pi$ up to $w$ is a shortest path from $u$ to $w$, because a
cheaper prefix would give a cheaper $\pi$.

\begin{definition}[Single-source shortest-path problem]
\label{def:ch03-sssp}
Given $G=(V,E,c)$ with $c\ge0$ and a source $s\in V$, compute $\delta(s,v)$ for every $v\in V$
together with a \textbf{shortest-path tree}\index{shortest-path tree}: a parent pointer
$\mathrm{parent}(v)$ for every reachable $v\neq s$ such that following the pointers from $v$
back to $s$ traces a shortest path. In the \textbf{single-pair} variant a target $t$ is also
given, and only $\delta(s,t)$ and one shortest path from $s$ to $t$ are required.
\end{definition}

During the search every vertex carries a \textbf{tentative cost}\index{tentative cost}
$\gcost(v)$, the cost of the best path to $v$ found so far, so that $\gcost(v)\ge\delta(s,v)$ at
all times. The open vertices form the priority queue \Open and the settled vertices the set
\Closed, exactly as in the best-first skeleton of \cref{ch:ch02}. In later chapters $\gcost$ is
joined by a heuristic $\hcost$; in this chapter $\hcost=0$ and the queue is ordered by $\gcost$
alone.

% ---------------------------------------------------------------------
\section{The algorithm}
\label{sec:ch03-algorithm}

\Cref{alg:ch03-dijkstra} is Dijkstra's algorithm in the form that every search algorithm of this
book reuses. A binary min-heap holds pairs $(\gcost,v)$, and instead of decreasing the key of a
vertex that is already in the heap, we push a second entry and ignore the older one when it
surfaces. This is the \textbf{lazy deletion}\index{priority queue!lazy deletion}\index{lazy
deletion} idiom of \cref{ch:ch02}. An optional target $t$ turns the algorithm into uniform-cost
search.

\begin{algorithm}[htb]
\caption{Dijkstra's algorithm with a binary heap and lazy deletion. Without a target it computes
$\delta(s,v)$ for every vertex; with a target $t$ it stops as soon as $t$ is settled
(uniform-cost search).}
\label{alg:ch03-dijkstra}
\KwIn{graph $G=(V,E,c)$ with $c\ge0$, given by a successor function $\Succ$; source $s$;
optional target $t$}
\KwOut{tentative costs $\gcost$, exact for every settled vertex; parent pointers}
\Fn{\DijkstraSearch{$G$, $s$, $t$}}{
  $\gcost(v)\gets\infty$ for all $v\in V$; $\gcost(s)\gets0$; $\mathrm{parent}(s)\gets$
\KwNil\label{alg:ch03-dijkstra:init}\;
  \Open $\gets$ min-heap holding the single entry $(0,s)$; \Closed
$\gets\emptyset$\label{alg:ch03-dijkstra:queue}\;
  \While{\Open $\neq\emptyset$}{
    $(g,u)\gets$ pop the entry with the smallest $g$ from \Open\label{alg:ch03-dijkstra:pop}\;
    \If{$u\in$ \Closed\label{alg:ch03-dijkstra:stale}}{
      \KwContinue\tcp*{stale entry: $u$ was settled earlier with a smaller key}
    }
    add $u$ to \Closed\label{alg:ch03-dijkstra:settle}\tcp*{settle $u$: from now on
$\gcost(u)=\delta(s,u)$}
    \If{$u=t$\label{alg:ch03-dijkstra:exit}}{
      \KwBreak\tcp*{early exit: uniform-cost search}
    }
    \ForEach{$v\in\Succ(u)$ with $v\notin$ \Closed\label{alg:ch03-dijkstra:succ}}{
      \If{$g+c(u,v)<\gcost(v)$\label{alg:ch03-dijkstra:test}}{
        $\gcost(v)\gets g+c(u,v)$; $\mathrm{parent}(v)\gets
u$\label{alg:ch03-dijkstra:relax}\tcp*{relaxation}
        push $(\gcost(v),v)$ onto \Open\label{alg:ch03-dijkstra:push}\tcp*{lazy insertion, no
decrease-key}
      }
    }
  }
  \KwRet{$\gcost$, $\mathrm{parent}$}\label{alg:ch03-dijkstra:ret}\;
}
\end{algorithm}

\paragraph{Walkthrough.} Lines~\ref{alg:ch03-dijkstra:init}--\ref{alg:ch03-dijkstra:queue}
initialise: only the source has a finite cost and only the source is on the frontier. In an
implementation the assignment $\gcost(v)\gets\infty$ for all $v$ is never executed; a dictionary
that answers $\infty$ for a missing key does the same job without touching $V$, which matters
when $V$ is a large grid or the space-time graph of \cref{ch:ch08}.

Line~\ref{alg:ch03-dijkstra:pop} removes the entry with the smallest key; with a binary heap
this costs $\bigO{\log\abs{\text{\Open}}}$. Line~\ref{alg:ch03-dijkstra:stale} is the
lazy-deletion test. A vertex that was pushed several times sits in the heap several times, and
every entry of it after the first one to be popped is \textbf{stale}\index{stale entry}: its key
is at least the key with which the vertex was settled, and it is simply discarded.
Line~\ref{alg:ch03-dijkstra:settle} settles $u$; this is the moment at which
\cref{thm:ch03-invariant} guarantees $\gcost(u)=\delta(s,u)$.

Line~\ref{alg:ch03-dijkstra:exit} is the early exit\index{Dijkstra's algorithm!early exit}. Note
where it sits: the target is tested when it is \emph{popped}, not when it is first discovered on
line~\ref{alg:ch03-dijkstra:relax}. A target that has just been discovered may still be reached
more cheaply through a vertex that is waiting in the queue; a target that has been popped
cannot. \Cref{sec:ch03-example} shows both cases on a concrete graph.

Lines~\ref{alg:ch03-dijkstra:succ}--\ref{alg:ch03-dijkstra:push} \emph{expand} $u$: they
generate its successors and relax every edge. The strict inequality on
line~\ref{alg:ch03-dijkstra:test} keeps the first-found path when two paths cost the same; this
is a tie-breaking rule, and \cref{sec:ch03-implementation} discusses its consequences. Skipping
settled successors on line~\ref{alg:ch03-dijkstra:succ} is an optimisation, not a necessity:
with $c\ge0$ the test on line~\ref{alg:ch03-dijkstra:test} would fail for them anyway. The word
\textbf{expansion}\index{expansion} names the unit in which we count the work of a search: one
expansion is one non-stale pop together with the generation of its successors. The number of
expansions is what \cref{fig:ch03-expansions} and the comparisons of \cref{ch:ch04} measure.

\paragraph{The result.} After the loop, $\gcost(v)=\delta(s,v)$ for every settled vertex, and
$\gcost(v)=\infty$ for every vertex the wave never reached because there is no path from $s$.
Without a target every reachable vertex is settled. With a target, only the vertices with
$\delta(s,v)\le\delta(s,t)$ have been settled, and the tentative costs of the vertices left in
the queue are upper bounds, not answers. The shortest path to a settled vertex $v$ is read off
the parent pointers\index{parent pointer}: start at $v$, move to $\mathrm{parent}(v)$, and
repeat until the source is reached; reversing the list gives the path from $s$ to $v$
(\textbf{path reconstruction}\index{path reconstruction}). If $t$ is unreachable, the queue runs
empty before $t$ is popped and the search reports failure, which for a drone means that the goal
is sealed off by obstacles or no-fly zones and the map, not the planner, has to change.

% ---------------------------------------------------------------------
\section{The algorithm: detecting conflicts}
\label{sec:ch07-detection}

Given $k$ time-indexed paths, we want the \emph{first} vertex or swapping conflict, or the
verdict that there is none. ``First'' means earliest in time; within one time step, a vertex
conflict in the snapshot at $t$ comes before a swapping conflict during the step from $t$ to
$t+1$, which comes before the snapshot at $t+1$. Returning the first conflict, rather than any
conflict, is what \cbs needs (\cref{ch:ch09}), and it makes the output deterministic.

Two small ingredients handle paths of different lengths. The function \MapfPos
(line~\ref{alg:ch07-first-conflict:pad} of \cref{alg:ch07-first-conflict}) returns the entry
$\pi_i[t]$ for $t\le T_i$ and the goal $g_i$ afterwards; this is the stay-at-target convention
turned into code, and it is the only place where the convention enters. The
\textbf{horizon}\index{horizon of a plan} $T=\max_i T_i$
(line~\ref{alg:ch07-first-conflict:horizon}) is the last time step that needs checking:
after $T$ every agent is parked, so no new vertex conflict can appear and no swap can happen
(\cref{thm:ch07-detection}).

\begin{algorithm}[htb]
\caption{Pairwise conflict detection}
\label{alg:ch07-first-conflict}
\KwIn{time-indexed paths $\pi_1,\dots,\pi_k$ in canonical form, with costs $T_1,\dots,T_k$}
\KwOut{the earliest vertex or swapping conflict, or \KwNil if the plan is conflict-free}
\Fn{\MapfPos{$\pi_i$, $t$}}{
  \lIf{$t\le T_i$}{\KwRet{$\pi_i[t]$}}
  \KwRet{$g_i$}\tcp*{stay at target: pad with the goal}\label{alg:ch07-first-conflict:pad}
}
\Fn{\MapfFirstConflict{$\pi_1,\dots,\pi_k$}}{
  $T\gets\max_i T_i$\label{alg:ch07-first-conflict:horizon}\;
  \For{$t\gets 0$ \KwTo $T$}{
    \ForEach{pair $i<j$\label{alg:ch07-first-conflict:vertex}}{
      $v\gets$ \MapfPos{$\pi_i$, $t$}\;
      \lIf{$v=$ \MapfPos{$\pi_j$, $t$}}{\KwRet{vertex conflict $\langle a_i,a_j,v,t\rangle$}}
    }
    \If{$t<T$}{
      \ForEach{pair $i<j$\label{alg:ch07-first-conflict:swap}}{
        $u_i,v_i\gets$ \MapfPos{$\pi_i$, $t$}, \MapfPos{$\pi_i$, $t+1$};\quad
        $u_j,v_j\gets$ \MapfPos{$\pi_j$, $t$}, \MapfPos{$\pi_j$, $t+1$}\;
        \lIf{$u_i\neq v_i$ \KwAnd $u_i=v_j$ \KwAnd $v_i=u_j$}{\KwRet{swapping conflict $\langle
a_i,a_j,u_i,v_i,t\rangle$}}
      }
    }
  }
  \KwRet{\KwNil}\;
}
\end{algorithm}

\paragraph{Walkthrough of \cref{alg:ch07-first-conflict}.}
The outer loop visits the columns of the plan table in order. At column $t$,
line~\ref{alg:ch07-first-conflict:vertex} compares every pair of agents and reports the
first pair that shares a vertex. If the column is clean and it is not the last one,
line~\ref{alg:ch07-first-conflict:swap} looks at the step from $t$ to $t+1$: for each pair
it reads the two moves $u_i\to v_i$ and $u_j\to v_j$ and reports a swap if they are the
reverse of each other. The test $u_i\neq v_i$ excludes the case where both agents wait on
the same vertex, which is a vertex conflict already reported at $t$. The order of the pairs
in the two inner loops (increasing $i$, then increasing $j$) fixes which conflict is reported
when several exist at the same time; the \python version follows the same order. There are
$\binom{k}{2}$ pairs per column and $T+1$ columns, each test costs constant time, so the
running time is $\bigO{k^2\,T}$\index{conflict detection!pairwise}. The memory beyond the
input is constant.

\paragraph{The hash-map version.}
A vertex conflict at column $t$ is nothing but a \emph{duplicate entry} in that column, and
duplicates in a list of $k$ items are found in expected $\bigO{k}$ time with a hash map.
\Cref{alg:ch07-hashed} keeps a map from vertices to the agent that occupies them
(line~\ref{alg:ch07-hashed:occ}); an agent whose vertex is already in the map conflicts with
the agent stored there. Swaps are found the same way with a map from directed edges to
agents (line~\ref{alg:ch07-hashed:mov}): agent $a_i$ moving $u\to v$ swaps with whoever is
stored under the reverse edge $(v,u)$. Each column now costs $\bigO{k}$ expected time, so the
whole scan costs $\bigO{k\,T}$\index{conflict detection!hash map} expected time and
$\bigO{k}$ extra memory. The saving matters: a planner that validates thousands of candidate
plans for a hundred agents cannot afford $\binom{100}{2}$ comparisons per time step.

\begin{algorithm}[htb]
\caption{Conflict detection with hash maps}
\label{alg:ch07-hashed}
\KwIn{time-indexed paths $\pi_1,\dots,\pi_k$ in canonical form}
\KwOut{the earliest vertex or swapping conflict, or \KwNil}
\Fn{\MapfFirstConflictHashed{$\pi_1,\dots,\pi_k$}}{
  $T\gets\max_i T_i$\;
  \For{$t\gets 0$ \KwTo $T$}{
    $\mathit{occupied}\gets$ empty map from vertices to agents\label{alg:ch07-hashed:occ}\;
    \For{$i\gets 1$ \KwTo $k$}{
      $v\gets$ \MapfPos{$\pi_i$, $t$}\;
      \lIf{$v\in\mathit{occupied}$}{\KwRet{vertex conflict $\langle
a_{\mathit{occupied}[v]},a_i,v,t\rangle$}}
      $\mathit{occupied}[v]\gets i$\;
    }
    \lIf{$t=T$}{\KwBreak}
    $\mathit{moving}\gets$ empty map from directed edges $(u,v)$ to
agents\label{alg:ch07-hashed:mov}\;
    \For{$i\gets 1$ \KwTo $k$}{
      $u,v\gets$ \MapfPos{$\pi_i$, $t$}, \MapfPos{$\pi_i$, $t+1$}\;
      \lIf{$u=v$}{\KwContinue}
      \lIf{$(v,u)\in\mathit{moving}$}{\KwRet{swapping conflict $\langle
a_{\mathit{moving}[(v,u)]},a_i,v,u,t\rangle$}}
      $\mathit{moving}[(u,v)]\gets i$\;
    }
  }
  \KwRet{\KwNil}\;
}
\end{algorithm}

One detail makes the two algorithms interchangeable. When a column holds several
conflicts, \cref{alg:ch07-hashed} as written returns the one it meets first while scanning
$i=1,\dots,k$, which need not be the lexicographically smallest pair $(i,j)$ that
\cref{alg:ch07-first-conflict} returns. The \python implementation
(\cref{lst:ch07-hashed}) finishes the column and keeps the smallest pair, so that both
detectors return exactly the same conflict; the self-test checks this on sixty random plans.

\paragraph{Validating a plan.}
A solver's output, a plan read from a file, or a plan repaired after a replan should be
\emph{validated} before anybody trusts it. \Cref{alg:ch07-validate} checks that each path is a
legal time-indexed path from its start to its goal---right endpoints, every vertex in $V$,
every step a wait or a move along an edge---and then runs conflict detection. Path legality
costs $\bigO{\sum_i T_i}$, so the whole check is dominated by detection. The listing in
\cref{sec:ch07-implementation} returns \emph{all} problems rather than the first one, which
is more useful when debugging a solver.

\begin{algorithm}[htb]
\caption{Plan validator}
\label{alg:ch07-validate}
\KwIn{instance $\langle G,s,g\rangle$, plan $\pi=(\pi_1,\dots,\pi_k)$}
\KwOut{\KwTrue if $\pi$ is a valid solution, otherwise the list of problems found}
\Fn{\MapfValidatePlan{$G$, $s$, $g$, $\pi$}}{
  $\mathit{problems}\gets$ empty list\;
  \For{$i\gets 1$ \KwTo $k$}{
    \lIf{$\pi_i[0]\neq s_i$ \KwOr $\pi_i[T_i]\neq g_i$}{append ``wrong endpoints of $a_i$''}
    \For{$t\gets 0$ \KwTo $T_i$}{
      \lIf{$\pi_i[t]\notin V$}{append ``$a_i$ leaves the graph at $t$''}
      \lIf{$t<T_i$ \KwAnd $\pi_i[t+1]\neq\pi_i[t]$ \KwAnd $\set{\pi_i[t],\pi_i[t+1]}\notin
E$}{append ``illegal step of $a_i$ at $t$''}
    }
  }
  \If{$\mathit{problems}$ is empty}{
    $c\gets$ \MapfFirstConflict{$\pi_1,\dots,\pi_k$}\;
    \lIf{$c\neq$ \KwNil}{append $c$}
  }
  \lIf{$\mathit{problems}$ is empty}{\KwRet{\KwTrue}}
  \KwRet{$\mathit{problems}$}\;
}
\end{algorithm}

% ---------------------------------------------------------------------
\section{A worked example}
\label{sec:ch07-example}

\begin{example}[Three agents on a $4\times4$ grid]
\label{ex:ch07-worked}
\Cref{fig:ch07-example-instance} shows a $4\times4$ grid with obstacles at $(1,0)$ and
$(1,2)$ and three agents: $a_1$ from $(0,0)$ to $(0,3)$, $a_2$ from $(0,3)$ to $(0,0)$, and
$a_3$ from the pocket $(1,1)$ to $(0,1)$. Cell $(0,0)$ is a dead end whose only neighbour is
$(0,1)$, and $(0,2)$ connects only to $(0,1)$ and $(0,3)$, so every path of $a_1$ and of
$a_2$ must pass through $(0,1)$ and $(0,2)$, and $(0,1)$ is also the goal of $a_3$. The
instance is the function \code{worked\_example()} of \code{ch07\_mapf.py}; in the code the
agents are numbered $0,1,2$.
\end{example}

\begin{figure}[tb]
  \centering
  \inputfigure{ch07/example-instance}
  \caption{The instance of \cref{ex:ch07-worked} with the independent shortest paths
  (the naive plan). All three agents need cell $(0,1)$ within two time steps, and $a_1$
  and $a_2$ need the corridor $(0,1)$--$(0,2)$ in opposite directions.}
  \label{fig:ch07-example-instance}
\end{figure}

\paragraph{The naive plan.}
Planning each agent alone with breadth-first search (\code{naive\_plan}) gives the paths of
\cref{tab:ch07-naive-paths}. Their costs are $3$, $3$ and $1$, so if the plan were valid
it would have $\sumcost=7$ and $\makespan=3$; by \cref{prop:ch07-bounds} these are lower
bounds for every valid solution. The path of $a_3$ has only two entries, so its column
entries for $t\ge 2$ are the goal $(0,1)$ by stay-at-target; the table shows them in grey.

\begin{table}[tb]
  \centering\small
  \caption{Naive plan of \cref{ex:ch07-worked}: independent shortest paths, written as the
  plan table. Grey entries are the stay-at-target padding of $a_3$.}
  \label{tab:ch07-naive-paths}
  \begin{tabular}{@{}lccccc@{}}
  \toprule
  Agent & $t=0$ & $t=1$ & $t=2$ & $t=3$ & $T_i$\\
  \midrule
  $a_1$ & $(0,0)$ & $(0,1)$ & $(0,2)$ & $(0,3)$ & 3\\
  $a_2$ & $(0,3)$ & $(0,2)$ & $(0,1)$ & $(0,0)$ & 3\\
  $a_3$ & $(1,1)$ & $(0,1)$ & \textcolor{sbGray}{$(0,1)$} & \textcolor{sbGray}{$(0,1)$} & 1\\
  \bottomrule
  \end{tabular}
\end{table}

\paragraph{Detection by hand.}
\Cref{tab:ch07-detection-trace} runs \cref{alg:ch07-first-conflict} on the naive plan with
$T=3$. The first column, $t=0$, holds three distinct cells. The step from $0$ to $1$ has
no pair of reversed moves. Column $t=1$ holds $(0,1)$ twice, for $a_1$ and $a_3$, so the
algorithm stops and returns the vertex conflict $\langle a_1,a_3,(0,1),1\rangle$. The
table continues the scan to list every conflict, as \code{all\_conflicts} does: the step
from $1$ to $2$ has $a_1$ moving $(0,1)\to(0,2)$ and $a_2$ moving $(0,2)\to(0,1)$, a
swapping conflict; column $t=2$ has $a_2$ arriving at $(0,1)$ where $a_3$ is parked, a
vertex conflict that only exists because of stay-at-target; the rest is clean. Running
\code{python3 ch07\_mapf.py} prints exactly these three conflicts.

\begin{table}[tb]
  \centering\small
  \caption{Trace of \cref{alg:ch07-first-conflict} on \cref{tab:ch07-naive-paths}. The
  algorithm returns at the third check; the remaining rows show what a full scan finds.}
  \label{tab:ch07-detection-trace}
  \begin{tabular}{@{}lL{5.6cm}L{5.9cm}@{}}
  \toprule
  Check & Column entries or moves & Result\\
  \midrule
  vertex, $t=0$ & $(0,0),\ (0,3),\ (1,1)$ & all distinct\\
  swap, $0\to1$ & $a_1\!:(0,0)\to(0,1)$; $a_2\!:(0,3)\to(0,2)$; $a_3\!:(1,1)\to(0,1)$ & no
reversed pair\\
  vertex, $t=1$ & $(0,1),\ (0,2),\ (0,1)$ & \textbf{return} $\langle a_1,a_3,(0,1),1\rangle$\\
  \midrule
  swap, $1\to2$ & $a_1\!:(0,1)\to(0,2)$; $a_2\!:(0,2)\to(0,1)$; $a_3$ waits & swap $\langle
a_1,a_2,(0,1),(0,2),1\rangle$\\
  vertex, $t=2$ & $(0,2),\ (0,1),\ (0,1)$ & vertex $\langle a_2,a_3,(0,1),2\rangle$ (parked
$a_3$)\\
  swap, $2\to3$ & $a_1\!:(0,2)\to(0,3)$; $a_2\!:(0,1)\to(0,0)$; $a_3$ waits & no reversed pair\\
  vertex, $t=3$ & $(0,3),\ (0,0),\ (0,1)$ & all distinct\\
  \bottomrule
  \end{tabular}
\end{table}

\paragraph{A conflict-free plan.}
\Cref{tab:ch07-solution-paths} and \cref{fig:ch07-example-solution} show a valid solution.
Agent $a_3$ first steps \emph{down} to $(2,1)$ to free the pocket; $a_1$ enters the pocket
at $t=2$ so that $a_2$ can pass along the top row and reach its dead-end goal at $t=3$;
$a_1$ then comes back up and continues, and $a_3$ follows it into the pocket and up to its
goal. Run the detector on the table to confirm that no column has a duplicate and no step
has a reversed pair; note in particular that the step from $2$ to $3$ has $a_1$ moving
$(1,1)\to(0,1)$ while $a_2$ moves $(0,1)\to(0,0)$ and $a_3$ moves $(2,1)\to(1,1)$---two
following conflicts, both allowed. The costs are $T_1=5$, $T_2=3$, $T_3=4$, so
\[
  \sumcost = 5+3+4 = 12, \qquad \makespan = \max(5,3,4) = 5 .
\]
The brute-force search of the code reports $12$ as the optimal sum of costs and $5$ as the
optimal makespan, so this plan is optimal for both objectives. Compare with the lower
bounds $7$ and $3$: the difference of five time steps is the price of coupling. All three
agents need one cell, and the only way to order them through it is to have two of them
step aside.

\begin{table}[tb]
  \centering\small
  \caption{A valid solution of \cref{ex:ch07-worked}, optimal for both objectives.
  Grey entries are stay-at-target padding; $\sumcost=12$, $\makespan=5$.}
  \label{tab:ch07-solution-paths}
  \begin{tabular}{@{}lccccccc@{}}
  \toprule
  Agent & $t=0$ & $t=1$ & $t=2$ & $t=3$ & $t=4$ & $t=5$ & $T_i$\\
  \midrule
  $a_1$ & $(0,0)$ & $(0,1)$ & $(1,1)$ & $(0,1)$ & $(0,2)$ & $(0,3)$ & 5\\
  $a_2$ & $(0,3)$ & $(0,2)$ & $(0,1)$ & $(0,0)$ & \textcolor{sbGray}{$(0,0)$} &
\textcolor{sbGray}{$(0,0)$} & 3\\
  $a_3$ & $(1,1)$ & $(2,1)$ & $(2,1)$ & $(1,1)$ & $(0,1)$ & \textcolor{sbGray}{$(0,1)$} & 4\\
  \bottomrule
  \end{tabular}
\end{table}

\begin{figure}[tb]
  \centering
  \inputfigure{ch07/example-solution}
  \caption{The solution of \cref{tab:ch07-solution-paths} as six snapshots; arrows show
  the move of the coming step. Notice how $a_1$ dips into the pocket at $t=2$ to let $a_2$
  pass, and how $a_3$ vacates the pocket beforehand and refills it afterwards.}
  \label{fig:ch07-example-solution}
\end{figure}

% ---------------------------------------------------------------------
\section{Properties: detection, bounds and complexity}
\label{sec:ch07-properties}

\begin{theorem}[Correctness of conflict detection]
\label{thm:ch07-detection}
Let $\pi_1,\dots,\pi_k$ be time-indexed paths in canonical form and $T=\max_i T_i$.
\Cref{alg:ch07-first-conflict,alg:ch07-hashed} return \KwNil if and only if the plan has no
vertex and no swapping conflict at any time $t\ge0$ under stay-at-target; otherwise they
return the earliest such conflict. \Cref{alg:ch07-first-conflict} runs in $\bigO{k^2 T}$
time; \cref{alg:ch07-hashed} runs in $\bigO{kT}$ expected time.
\end{theorem}
\begin{proof}
By construction \MapfPos{$\pi_i$,$t$} equals the vertex of $a_i$ at time $t$ under
stay-at-target for every $t\ge0$, so the tests at columns $0,\dots,T$ and at the steps
between them are exactly the conditions of
\cref{def:ch07-vertex-conflict,def:ch07-swapping-conflict}. It remains to show that nothing
is missed after $T$. For $t\ge T$ every agent is at its goal and stays there, so
\MapfPos{$\pi_i$,$t$}$=$\MapfPos{$\pi_i$,$T$}. A vertex conflict at some $t>T$ is
therefore also a vertex conflict at $T$, which the scan checks. A swapping conflict during a
step $t\to t+1$ with $t\ge T$ needs two moving agents, but nobody moves after $T$. Hence
the scan up to $T$ finds a conflict if and only if one exists, and because it visits the
columns and the steps in chronological order, the conflict returned is the earliest. The
running times were derived in \cref{sec:ch07-detection}.
\end{proof}

\begin{proposition}[Bounds and relations between the objectives]
\label{prop:ch07-bounds}
Write $d(u,v)$ for the shortest-path distance in $G$. For every valid solution $\pi$ of
$\langle G,s,g\rangle$:
\begin{enumerate}[label=(\alph*)]
  \item $\sumcost(\pi)\ge\sum_i d(s_i,g_i)$ and $\makespan(\pi)\ge\max_i d(s_i,g_i)$;
  \item $\makespan(\pi)\le\sumcost(\pi)\le k\cdot\makespan(\pi)$;
  \item if the plan of independent shortest paths (the \textbf{naive plan}\index{naive plan})
        is conflict-free, it is optimal for both objectives.
\end{enumerate}
\end{proposition}
\begin{proof}
(a) $T_i$ counts all actions of $a_i$ before its final arrival, so $T_i\ge$ number of
moves $\ge d(s_i,g_i)$; sum or maximise over $i$. (b) The maximum of $k$ non-negative
numbers is at most their sum, and their sum is at most $k$ times their maximum. (c) The naive
plan attains the lower bounds of (a) and is valid, so nothing valid is cheaper.
\end{proof}

The bounds in (a) are the starting point of every optimal solver: \cbs begins with the
naive plan and only pays for the conflicts it has to resolve (\cref{ch:ch09}), and the
sum of the single-agent distances is the natural admissible heuristic for the joint search
of \mstar (\cref{ch:ch11}). \Cref{ex:ch07-objectives} shows that (b) is all one can say in
general: the two objectives can be optimised by different plans.

\begin{theorem}[Optimal MAPF is NP-hard]
\label{thm:ch07-nphard}
Finding a valid solution of a classical MAPF instance with minimum sum of costs, or with
minimum makespan, is NP-hard \cite{yu2013structure,surynek2010optimization}. Consequently,
unless $\mathrm{P}=\mathrm{NP}$, no algorithm solves every instance optimally in time
polynomial in $\abs{V}$ and $k$.\index{NP-hardness!optimal MAPF}
\end{theorem}
\begin{proof}[Proof sketch]
Both results are reductions from Boolean satisfiability. Surynek~\cite{surynek2010optimization}
showed that deciding whether a makespan below a given bound is achievable is NP-hard. Yu and
LaValle~\cite{yu2013structure} built, from a 3-SAT formula, a graph and a set of agents such
that each variable agent must commit to one of two equally short routes (``true'' or
``false'') and each clause gadget lets the agents through without delay only if the route of
at least one of its literals is chosen; the resulting instance has a solution matching the
single-agent lower bound if and only if the formula is satisfiable, and the construction
works for the makespan and for the sum of costs. The models used in the two papers differ
in minor motion rules, but the hardness carries over to the classical setting used here
(see \cite{stern2019mapf} for the discussion). Ma and colleagues later strengthened the
makespan result to an inapproximability bound: even a solution within a factor smaller than
$4/3$ of the optimal makespan is NP-hard to find, as summarised in \cite{stern2019mapf}.
\end{proof}

\paragraph{Where the hardness comes from.}
A single agent lives in a state space of size $\abs{V}$. The $k$ agents together live in the
space of configurations, of size up to $\abs{V}^k$, and a solver that treats the swarm as one
big agent explores that space---this is the joint-state blow-up that \cref{ch:ch11} starts
from. What the reduction shows is that the blow-up cannot be avoided in general: at each
bottleneck somebody has to yield, the choice of who yields is essentially binary, and the
choices are \emph{coupled}\index{coupling of agents}---letting $a_1$ pass first pushes $a_2$
into the way of $a_3$, and so on. In the worst case the number of consistent combinations
of these choices is exponential, which is exactly what a constraint tree (\cref{ch:ch09}) or a
priority ordering (\cref{ch:ch08}) has to search through. In practice most agents never
interact, and the successful solvers exploit that; \cref{fig:ch07-conflict-growth} shows how
quickly interactions nevertheless multiply.

\begin{figure}[tb]
  \centering
  \inputfigure{ch07/conflict-growth}
  \caption{Conflicts of the naive plan on random $16\times16$ grids with $20\%$ obstacles,
  averaged over $60$ seeded instances per $k$ (generated by
  \code{code/figures/gen\_ch07\_conflict\_growth.py}). Left: the number of vertex and
  swapping conflicts grows roughly quadratically with the number of agents (about $5$ at
  $k=10$, $21$ at $k=20$, $97$ at $k=40$). Right: beyond a dozen agents the naive plan is
  never conflict-free.}
  \label{fig:ch07-conflict-growth}
\end{figure}

\begin{proposition}[Feasibility is polynomial]
\label{prop:ch07-feasibility}
Deciding whether a classical MAPF instance is solvable, and constructing a valid solution
when it is, can be done in polynomial time. Kornhauser, Miller and
Spirakis~\cite{kornhauser1984pebble} proved this for \emph{pebble motion on graphs}, in which
one agent at a time moves into an empty neighbouring vertex; every such sequential move is a
classical MAPF step in which all other agents wait, so every solvable pebble-motion instance
is a solvable MAPF instance. Push-and-Rotate \cite{dewilde2014push} is a practical
polynomial-time algorithm that solves every solvable instance on a connected graph with at
least two unoccupied vertices; it is presented in
\cref{ch:ch11}.\index{feasibility!MAPF}\index{Push-and-Rotate}
\end{proposition}

The two results together are the central fact of Part~III: \emph{finding some plan is easy,
finding a good plan is hard}. Push-and-Rotate finds valid but long plans; prioritized
planning (\cref{ch:ch08}) finds short plans quickly but sometimes none at all; \cbs
(\cref{ch:ch09}) finds optimal plans at exponential worst-case cost; \ecbs (\cref{ch:ch10})
buys speed with a bounded loss of quality. The next section lays these out.

% ---------------------------------------------------------------------
\section{Implementation notes}
\label{sec:ch03-implementation}

\Cref{lst:ch03-core} is the function \code{dijkstra} of \code{code/ch03\_dijkstra.py}, a direct
transcription of \cref{alg:ch03-dijkstra}. The remarks below explain the choices in it;
\cref{lst:ch03-helpers} adds path reconstruction and the backward run.

\paragraph{The heap.} \code{heapq}\index{heapq@\code{heapq}} is a binary min-heap stored in a
list; it offers \code{heappush} and \code{heappop} but no decrease-key, which is why the listing
pushes duplicates. Do not ``fix'' this by searching the list for the old entry: that costs
$\bigO{\abs{\text{\Open}}}$ per relaxation and destroys the bound of \cref{thm:ch03-complexity}.
Entries are tuples \code{(g, node)}; \python compares tuples lexicographically, so ties in
\code{g} are broken by comparing the nodes themselves. This works for strings and for
\code{(row, col)} tuples and makes the search deterministic, which is what you want when you
compare runs or reproduce a trace. If your nodes are not comparable, insert a running counter as
the second element of the tuple, as \code{LazyPQ} in \cref{ch:ch02} does.

\paragraph{Costs and parents.} \code{dist} is a dictionary: a missing key means $\infty$, so
\code{dist.get(v, INF)} is the test of line~\ref{alg:ch03-dijkstra:test} and no array of size
$\abs{V}$ is ever allocated. \code{parent} maps every discovered node to its predecessor and the
source to \code{None}; \code{reconstruct\_path} follows the pointers back and reverses the list.
Because the parent is overwritten at every relaxation, the pointers always describe the path
that produced the current \code{dist}.

\paragraph{The closed set.} \code{closed} is a set, and the membership test on pop is the
lazy-deletion check. After the loop, \code{dist[v]} is exact if and only if \code{v in closed};
with an early exit, the tentative costs of open nodes are upper bounds only (recall $F$ in
\cref{ex:ch03-graph}). If your caller needs the guarantee, return \code{closed} as well, or
return only the settled entries.

\paragraph{Implicit graphs.} The function only ever evaluates \code{graph[u]}, which must yield
\code{(neighbour, cost)} pairs. A dictionary of lists does this for explicit graphs. The class
\code{GridGraph} in the same file does it for occupancy grids: it generates the 4 or 8
neighbours of a cell on the fly, with cost 1 or $\sqrt2$, forbids corner cutting as in
\cref{ch:ch02}, and multiplies the cost of a move by an optional factor attached to the cell
being entered. Its method \code{reversed()} flips the direction in which the factors are
applied, so that \cref{alg:ch03-backward} can run on a grid without ever building the reversed
graph explicitly.

\paragraph{Tie-breaking.}\index{tie-breaking} Among several shortest paths, the one returned
depends on two rules: the strict \code{<} on line~\ref{alg:ch03-dijkstra:test}, which keeps the
first path found at equal cost, and the order of the heap among equal keys. Both are
deterministic here, but neither is canonical: a different neighbour order yields a different,
equally optimal path. Planners that must be reproducible fix the neighbour order, as
\code{MOVES} in the code does. \astar in \cref{ch:ch04} adds a tie-breaking rule on $\gcost$
that changes the number of expansions substantially.

\paragraph{Floating-point costs.}\index{floating point!cost comparison} On 8-connected grids the
costs $1$ and $\sqrt2$ are added in different orders along different paths, so two paths of the
same true cost can end with values that differ in the last binary digit. Never test costs for
equality: compare with a tolerance, or scale the costs to integers (many grid planners use 10
for a straight move and 14 for a diagonal one). The self-test of the chapter code compares
against a brute-force enumeration with a tolerance of $10^{-9}$ and checks that $\hcost^*$ from
the backward run agrees with a forward search from every cell.

\begin{lstlisting}[caption={The core of the implementation (from \code{code/ch03\_dijkstra.py}): the main loop with \code{heapq}, a dictionary of tentative costs, a closed set tested on every pop, and lazy insertion instead of decrease-key.},label={lst:ch03-core}]
def dijkstra(graph, source, target=None):
    """Shortest paths from source; every edge cost must be >= 0.

    graph[u] yields the (v, cost) pairs of the out-edges of u; a dict of
    lists or a GridGraph both work.  Returns (dist, parent): dist[v] is the
    cost of the best path found to v and parent[v] its predecessor on that
    path (parent[source] is None).  Unreached nodes are absent (infinity).
    With a target the loop stops as soon as the target is settled
    (uniform-cost search); dist is then exact for settled nodes only.
    """
    dist = {source: 0.0}
    parent = {source: None}
    closed = set()
    queue = [(0.0, source)]            # min-heap of (g, node) entries
    while queue:
        g, u = heapq.heappop(queue)    # smallest g first, ties by node
        if u in closed:                # stale entry: u was settled earlier
            continue
        closed.add(u)                  # settle u: dist[u] == g is final
        if u == target:
            break                      # early exit
        for v, cost in graph[u]:
            if v in closed:
                continue
            g_new = g + cost
            if g_new < dist.get(v, INF):
                dist[v] = g_new        # relaxation
                parent[v] = u
                heapq.heappush(queue, (g_new, v))   # lazy insertion
    return dist, parent
\end{lstlisting}

\begin{lstlisting}[caption={Path reconstruction from the parent pointers and the backward run that produces the true-distance heuristic table (from \code{code/ch03\_dijkstra.py}).},label={lst:ch03-helpers}]
def reconstruct_path(parent, target):
    """Follow parent pointers back from target; None if unreachable."""
    if target not in parent:
        return None
    path = []
    node = target
    while node is not None:
        path.append(node)
        node = parent[node]
    path.reverse()
    return path


def backward_dijkstra_heuristic(graph, goal):
    """True-distance heuristic h*(v) = dist(v, goal) for every node v.

    One run of dijkstra() from the goal over the reversed graph.  The
    table is a perfect heuristic: admissible and consistent on the static
    map, and still admissible when wait actions or constraints are added
    (Chapters 4, 8 and 9).  Unreached nodes are absent: read as infinity.
    """
    if hasattr(graph, "reversed"):
        rev = graph.reversed()         # GridGraph knows how to flip itself
    else:
        rev = reverse_graph(graph)
    h_star, _ = dijkstra(rev, goal)
    return h_star
\end{lstlisting}

\begin{pitfall}[Testing for the goal at the wrong time]
The most common bug in a first implementation is to return as soon as the target is
\emph{discovered}, on line~\ref{alg:ch03-dijkstra:relax}, instead of when it is \emph{popped}.
In \cref{ex:ch03-graph} the vertex $F$ is discovered with cost 13, later improved to 12 and then
to 11. A search that stops at discovery would report 13, almost twenty percent too much, and the
``optimal'' path it returns would be $A,B,C,D,F$. Breadth-first search may test at discovery
because every edge costs the same; Dijkstra's algorithm and \astar may not. The same rule
applies to the early exit of every search in \cref{ch:ch04,ch:ch08,ch:ch09}.
\end{pitfall}

\begin{pitfall}[Negative costs and the constant-shift trap]
Dijkstra's algorithm is silently wrong when an edge has a negative cost
(\cref{ex:ch03-negative}): it does not crash, it returns a path that is not optimal. Negative
costs sneak into planners as rewards, for instance a bonus for flying through a cell that must
be surveyed. Adding a constant to all edge costs does not help, because paths with more edges
are penalised more and the optimal path changes (\cref{exr:ch03-shift}). Model rewards as
reduced non-negative costs, or use the Bellman--Ford algorithm. Zero costs are fine.
\end{pitfall}

\begin{pitfall}[Forgetting the stale check]
Without the test on line~\ref{alg:ch03-dijkstra:stale} every duplicate entry is expanded again.
The costs stay correct, because a relaxation from a stale entry can never improve anything, but
a vertex is then ``settled'' several times: the count of expansions is wrong, a visualisation of
the closed set shows cells being settled twice, and in the incremental algorithms of
\cref{ch:ch05}, where an expansion has side effects, the result is wrong. Test membership in the
closed set, or compare the popped key with the current \code{dist[u]}; either works, but one of
them must be there.
\end{pitfall}

% ---------------------------------------------------------------------
\section{Where this fits in the drone system}
\label{sec:ch03-drone}

\begin{dronebox}
\index{Dijkstra's algorithm!in the drone system}%
In the hybrid architecture of \cref{ch:ch24} Dijkstra's algorithm is part of the global planning
layer, and it plays two roles there. As the \emph{baseline planner} it routes a single drone
across a static occupancy grid or waypoint graph whose edge costs encode flight time, energy or
risk; its output, an optimal path with its cost, is the yardstick against which every faster or
more capable planner of the book is measured, and the experiments of \cref{ch:ch25} report path
lengths relative to it. As the \emph{heuristic factory} it runs backwards from every drone's
goal once, on the static map, and delivers the true-distance table $\hcost^*$ that \astar
(\cref{ch:ch04}), the space-time searches of prioritized planning (\cref{ch:ch08}) and the low
level of conflict-based search (\cref{ch:ch09}) consult at every expansion. The replanning layer
inherits the same backward wave in incremental form as \dstarlite (\cref{ch:ch05}). What
Dijkstra's algorithm does not do is react: it does not see other drones, and it does not see the
intruder that appears after take-off; those are handled by the multi-agent solver, by the
prediction layer and by the local avoidance layer. In the study plan (\cref{ch:appA}) this
chapter is the first half of Week~1; the second half is \astar.
\end{dronebox}

% ---------------------------------------------------------------------
\section{Summary}
\begin{summary}
\begin{itemize}
  \item Dijkstra's algorithm computes $\delta(s,v)$ for every vertex of a graph with
non-negative edge costs, together with a shortest-path tree, by settling vertices in
non-decreasing order of tentative cost: a wave spreading from the source.
  \item The settled-node invariant (\cref{thm:ch03-invariant}) says that a vertex's cost is
final when it is popped. Its proof uses $c\ge0$ exactly once, and a single negative edge breaks
it (\cref{ex:ch03-negative}).
  \item With a binary heap and lazy deletion the running time is
$\bigO{(\abs{V}+\abs{E})\log\abs{V}}$, the number of expansions grows with the area of the wave,
and the early exit (uniform-cost search) stops the wave at the radius of the target.
  \item The implementation needs a heap of $(\gcost,v)$ pairs, a dictionary of tentative costs,
a closed set that is tested on every pop, parent pointers for the path, and tolerance-based
comparisons of floating-point costs.
  \item One backward run from the goal on the reversed graph gives the true-distance heuristic
$\hcost^*(v)=\delta(v,t)$, which is consistent, exact on the static map, and still admissible in
space-time and under constraints: the heuristic used throughout \cref{ch:ch04,ch:ch08,ch:ch09}.
\end{itemize}
\end{summary}

\paragraph{Further reading.}
The original three-page paper is \cite{dijkstra1959note}; the array-scanning version in it is
still the right choice for dense graphs. \textcite{cormen2009clrs} give the standard textbook
treatment with the proof of the invariant, the Bellman--Ford algorithm for negative costs and
the priority-queue trade-offs, and \textcite{fredman1987fibonacci} introduce the Fibonacci heap.
\textcite{russell2020aima} present the algorithm as uniform-cost search alongside breadth-first
and bidirectional search, and \textcite{felner2011position} explains why that presentation is
the better one for implicit graphs. The step from Dijkstra's algorithm to \astar was taken by
\textcite{hart1968formal}; \textcite{pearl1984heuristics} is the classic reference on
heuristics, including the exact heuristic of \cref{sec:ch03-backward}.

% ---------------------------------------------------------------------
\section{Exercises}
\label{sec:ch03-exercises}

\begin{exercise}[\difficulty{1}]
\label{exr:ch03-trace}
In \cref{ex:ch03-graph} reduce the cost of the edge $B$--$E$ from 7 to 4 and trace
\cref{alg:ch03-dijkstra} from $A$ by hand in the format of \cref{tab:ch03-trace}. Give the final
distances, the shortest path to $G$, the number of pops and the number of stale entries. Check
your trace with \code{dijkstra\_trace} in the chapter code.
\end{exercise}

\begin{exercise}[\difficulty{1}]
\label{exr:ch03-early}
Suppose \cref{alg:ch03-dijkstra} returned as soon as the target is pushed on
line~\ref{alg:ch03-dijkstra:push} rather than when it is popped. Using \cref{tab:ch03-trace},
give the cost and path it would return for the target $F$, and for the target $G$. Which step of
the proof of \cref{thm:ch03-invariant} no longer applies to a vertex that has been discovered
but not settled?
\end{exercise}

\begin{exercise}[\difficulty{2}]
\label{exr:ch03-monotone}
Use \cref{thm:ch03-monotone,thm:ch03-invariant} to prove that uniform-cost search with target
$t$ settles every vertex $v$ with $\delta(s,v)<\delta(s,t)$, no vertex with
$\delta(s,v)>\delta(s,t)$, and that which vertices with $\delta(s,v)=\delta(s,t)$ it settles
depends only on tie-breaking. Then show that the tentative cost of every vertex left in the
queue at the exit is at least $\delta(s,t)$, and explain why it may nevertheless differ from
$\delta(s,v)$. Finally, show that edges of cost 0, including zero-cost cycles, do not endanger
termination.
\end{exercise}

\begin{exercise}[\difficulty{2}]
\label{exr:ch03-shift}
Let $G$ have some negative edge costs but no negative cycle, and let $G'$ be obtained by adding
the same constant $K>0$ to every edge cost so that all costs become non-negative. Show with a
graph of four vertices that a shortest path in $G'$ need not be a shortest path in $G$. Then
explain, in one sentence, what property of paths makes the shift harmless when all shortest
paths of interest have the same number of edges.
\end{exercise}

\begin{exercise}[\difficulty{2}]
\label{exr:ch03-grid8}
Repeat \cref{ex:ch03-grid} on the same grid with 8-connectivity: a diagonal move costs $\sqrt2$,
entering a hatched cell multiplies the cost of the move by 3, and corner cutting is forbidden.
Determine $\delta(S,T)$ and an optimal path by hand, then run \code{dijkstra\_trace} on
\code{GridGraph(cells, 8, cost)} and count the expansions and the stale entries. Why does this
instance produce stale entries when the 4-connected one did not?
\end{exercise}

\begin{exercise}[\difficulty{2}]
\label{exr:ch03-backward}
(a) Compute the true-distance heuristic $\hcost^*(v)=\delta(v,G)$ for every vertex of
\cref{ex:ch03-graph} by running \cref{alg:ch03-backward} by hand, and verify the consistency
inequality $\hcost^*(u)\le c(u,v)+\hcost^*(v)$ on every edge. (b) On the grid of
\cref{ex:ch03-grid}, explain why the forward table from $T$ and the backward table to $T$
differ, and name every cell where they differ. (c) A drone must reach $T$ on that grid but is
forbidden to enter the cell $(4,2)$ at time step 6 (a vertex constraint of the kind used in
\cref{ch:ch09}). Explain why the table of \cref{fig:ch03-backward-heuristic} is still an
admissible heuristic for the constrained space-time search, and why it is no longer exact.
\end{exercise}

\begin{exercise}[\difficulty{3} Coding]
\label{exr:ch03-coding}
This is the first half of the Week~1 exercise of the study plan (\cref{ch:appA}); \cref{ch:ch04}
completes it. Implement Dijkstra's algorithm on an occupancy grid from scratch, with 4- and
8-connectivity, a successor function that forbids corner cutting, lazy deletion, the early exit
and path reconstruction. (a) Visualise the closed set at several moments of a search on a
$30\times30$ grid with a few walls, as in \cref{fig:ch03-grid-closed}, and check that the
settled cells at any moment are exactly the cells whose cost is at most the last popped key. (b)
Count expansions against grid size for random grids with and without the early exit, and
reproduce the shape of \cref{fig:ch03-expansions}. (c) Add a function that returns the
true-distance table for a goal and check, on random grids, that $\hcost^*(v)$ equals the cost of
a forward search from $v$ for every free cell $v$. Keep the code: the next chapter turns it into
\astar by changing the key of the heap.
\end{exercise}

\begin{exercise}[\difficulty{3}]
\label{exr:ch03-bidirectional}
Implement bidirectional Dijkstra as described in \cref{sec:ch03-bidirectional}. (a) Construct a
small graph on which stopping at the first vertex reached by both waves returns a suboptimal
path, and show that the stopping rule with $\mu$ returns the optimal one. (b) Prove that the
rule is correct: when the two smallest keys sum to at least $\mu$, no path cheaper than $\mu$
exists. (c) On the random grids of \cref{fig:ch03-expansions}, measure the number of settled
cells of bidirectional search relative to uniform-cost search for the far-corner target, and
compare with the factor of one half predicted by the disc argument.
\end{exercise}

